\documentclass[a4paper, 12pt, numbers=noenddot, parskip=half]{scrartcl}

% Querverweise in das Dokument MehrheitsVerhältniswahl (muss vor hyperref geladen werden).
% Externe Labels werden mit dem Präfix MV- referenziert, z.B. \ref{MV-subsec:verfassungsaenderungen}.
% Voraussetzung: MehrheitsVerhältniswahl wurde vorher gebaut (dist/MehrheitsVerhältniswahl.aux).
\usepackage{xr-hyper}
% Kodierung und Sprache
\usepackage[utf8]{inputenc}
\usepackage[T1]{fontenc}
\usepackage[ngerman]{babel}

% Schrift: Helvetica (wie Arial in Google Docs)
\usepackage{helvet}
\renewcommand{\familydefault}{\sfdefault}

% Seitenformat
\usepackage[a4paper, margin=2.5cm]{geometry}

% Griechische Buchstaben im Text (\textchi etc., aus greek-fontenc)
\usepackage{textalpha}

% Mathematik
\usepackage[fleqn]{amsmath}
\setlength{\mathindent}{0pt}
\usepackage{amssymb}
\usepackage{mathtools}
\usepackage{amsthm}
\theoremstyle{definition}
\newtheorem{hilfssatz}{Hilfssatz}
\newtheorem{hauptsatz}{Hauptsatz}

% Diagramme
\usepackage{tikz}
\usetikzlibrary{calc, positioning, patterns}

% Farben für Balkendiagramme
\usepackage{xcolor}
\definecolor{parteiA}{gray}{0.25}
\definecolor{parteiB}{gray}{0.65}
\definecolor{parteiC}{RGB}{100,160,220}
\definecolor{parteiD}{RGB}{80,160,80}
\definecolor{parteiE}{RGB}{200,60,60}
\definecolor{parteiF}{RGB}{220,190,40}
\definecolor{parteiS}{RGB}{160,120,40}
\definecolor{diagrammHG}{RGB}{210,230,245}

% Tabellen
\usepackage{booktabs}
\usepackage{array}
\usepackage{tabularx}

% Listen
\usepackage{enumitem}

% Typografie
\usepackage{microtype}
\usepackage{csquotes}
\widowpenalty=10000
\clubpenalty=10000
\usepackage[ragged]{footmisc}
\usepackage{needspace}
\newcommand{\runinhead}[1]{\needspace{3\baselineskip}\textbf{#1}}
\newcommand{\runinheadul}[1]{\needspace{3\baselineskip}\underline{#1}}
% Mindestplatz vor Überschriften; direkt nach einer Überschrift entfällt er,
% damit die übergeordnete Überschrift nicht allein am Seitenende stehen bleibt
\makeatletter
\newcommand{\headingspace}[1]{\if@nobreak\else\needspace{#1\baselineskip}\fi}
\makeatother
\AddToHook{cmd/section/before}{\headingspace{8}}
\AddToHook{cmd/subsection/before}{\headingspace{7}}
\AddToHook{cmd/subsubsection/before}{\headingspace{5}}
\AddToHook{cmd/paragraph/before}{\headingspace{4}}
% \subsubsubsection ist eine Kopie von \paragraph (siehe unten) und braucht einen eigenen Hook
\AddToHook{cmd/subsubsubsection/before}{\headingspace{4}}

% Vierte Gliederungsebene: \subsubsubsection als eigenständige, nummerierte
% Überschrift unterhalb von \subsubsection (Alias für \paragraph)
\let\subsubsubsection\paragraph
\setcounter{secnumdepth}{4}
\setcounter{tocdepth}{4}
\RedeclareSectionCommand[
  beforeskip=-1.5ex plus -1ex minus -.2ex,
  afterskip=1em,
  font=\normalfont\bfseries,
  indent=0pt,
  tocraggedentrytext=true
]{paragraph}

% Grafiken (für Screenshots-Platzhalter im Anhang)
% float muss vor hyperref geladen werden, sonst entstehen doppelte Sprungziele (figure.N, table.N)
\usepackage{graphicx}
\usepackage{float}
\graphicspath{{assets/}{../assets/}}

% Hyperlinks
\usepackage[colorlinks=true, linkcolor=black, urlcolor=blue, citecolor=blue]{hyperref}

\usepackage{adjustbox}

% Hilfsmakro: Skalierung für Balkendiagramme (100\% = 8\,cm)
\newlength{\balkenhoehe}
\setlength{\balkenhoehe}{8cm}
\newcommand{\proztocm}[1]{\dimexpr#1\balkenhoehe/100\relax}

% Querverweise mit automatischer Bezeichnung, z.B. \cref{subsec:...} -> "Abschnitt 5.5"
% (muss nach hyperref und amsthm geladen werden, daher am Ende der Präambel)
\usepackage[ngerman,nameinlink]{cleveref}
% \subsubsubsection ist ein Alias für \paragraph, dafür kennt cleveref keinen Namen
\crefname{paragraph}{Abschnitt}{Abschnitte}
\Crefname{paragraph}{Abschnitt}{Abschnitte}

% PDF-Betrachter lesen den Dateinamen im Link als PDFDocEncoding (Latin-1), nicht als UTF-8.
% Deshalb steht das "ä" im Link als Byte ^^e4; ein direkt geschriebenes "ä" ergäbe "ä".
\externaldocument[MV-]{dist/MehrheitsVerhältniswahl}[MehrheitsVerh\string^^e4ltniswahl.pdf]
% Einheitlicher Verweis auf einen Abschnitt in MehrheitsVerhältniswahl (Label ohne MV-Präfix angeben)
\newcommand{\mvref}[1]{Dokument \enquote{Mehrheitssicherndes Verhältniswahlverfahren},
  Abschnitt~\ref{MV-#1} \enquote{\nameref{MV-#1}}}

\begin{document}

\begin{titlepage}
  \centering
  \vspace*{6cm}
  {\LARGE Vorschläge zur Änderung des Verfassungsgefüges\par}
  \vspace{2cm}
  {\normalsize Ein Diskussionsbeitrag\par}
  \vspace{1cm}
  {\normalsize Daniel Schwalbe\par}
  \vspace{1cm}
  {\normalsize Entwurf --- Oktober 2026\par}
  \vfill
\end{titlepage}

\tableofcontents
\newpage

\section{Notwendigkeit}

Änderungen im Verfassungsgefüge sind notwendig:
\begin{enumerate}[label=\alph*)]
    \item zur Härtung der Demokratie
    \begin{itemize}
        \item gegen Populismus
        \item gegen Kriminalität und Korruption in gewählten Ämtern
        \item gegen demokratisch gewählte Demokratiefeinde
        \item gegen Unfähigkeit der Spitzenpolitiker, vor allem des Bundeskanzlers
        \item gegen folgenschwere falsche Entscheidungen von Exekutive und Legislative
    \end{itemize}
    \item zur Anpassung an das vorgeschlagene neue Wahlsystem
\end{enumerate}

Die erforderlichen Änderungen sind so schwerwiegend, dass es kaum vorstellbar ist, sie innerhalb des normalen Verfassungsänderungsprozesses umzusetzen.
Sie bedürfen einer Legitimation durch das Volk selbst und nicht nur durch Abgeordnete des Volkes.
Meiner Meinung nach ist es jedoch ohnehin notwendig, dass sich das deutsche Volk in freier Entscheidung eine neue Verfassung gibt,
die die Lehren aus der Geschichte beider deutscher Staaten zieht und auch den Erfahrungen mit der Demokratie Rechnung trägt, die seit der Konstituierung der bundesdeutschen Verfassung 1949 weltweit gemacht wurden.
Im Rahmen eines solchen verfassunggebenden Prozesses können auch größere Änderungen im Verfassungsgefüge beschlossen werden.
Obwohl ich davon ausgehe, dass es eigentlich auf ein Neuschreiben der Verfassung hinauslaufen wird, werde ich jedoch im Folgenden nur die angedachten Änderungen zur aktuellen Verfassung beschreiben, damit ich keine neue Verfassung entwerfen muss.
Dazu würde mir auch die Qualifikation fehlen.
Schließlich bin ich Softwareentwickler und kein Politikwissenschaftler oder Verfassungsrechtler.

\section{Härtung der Demokratie gegen Populismus} \label{sec:haertung_der_demokratie}

Die Auseinandersetzung mit dem Populismus muss auf mehreren Ebenen erfolgen.
Wirklich besiegen kann man ihn nur, wenn man ihm die inhaltlichen Grundlagen entzieht.

Offenbar ist dies aber nicht so einfach, denn auch in der Demokratie gibt es Elitenbildung und, viel schlimmer noch, Abkopplung der Eliten von der Lebensrealität der Nicht-Eliten.
Die Existenz von Eliten an sich ist eher etwas Gutes: Wir alle wollen nicht von Menschen geführt werden, die dümmer sind als wir.
Eigentlich wollen wir auch nicht von Menschen geführt werden, die genauso klug sind wie wir.
Stattdessen wollen wir von den Besten unter uns geführt werden.

Problematisch wird es, wenn immer größere Teile der Bevölkerung den Eindruck gewinnen, dass
\begin{enumerate}[label=\alph*)]
    \item die Eliten in ihrer eigenen Welt leben, die mit der Realität nichts mehr zu tun hat, oder
    \item nur ihre eigenen Interessen verfolgen, die den Interessen des \enquote{Normalbürgers} entgegenlaufen, oder
    \item die Zeiten nicht richtig einschätzen können, oder
    \item einfach auch keine Lösung haben, oder
    \item die Eliten gar nicht wirklich klüger sind als sie, sondern nur privilegiert.
\end{enumerate}

Dem kann man eigentlich nur entgegenwirken, indem man neu Menschen von unten nach oben bringt.
Die Demokratie lässt dies im Gegensatz zur Monarchie ja auch ohne Revolution zu.

Allerdings macht sie es auch nicht unbedingt einfach, von unten nach oben aufzusteigen.
Mit wohlhabenden, gebildeten Eltern in der richtigen Partei hat man schon deutlich bessere Aufstiegschancen als ein Kind aus bildungs- und politikfernen Schichten.

Trotzdem ist es möglich, dass auch Menschen, die das Leben kennen, bis an die Spitze kommen.
Es mag aber dauern, bis es wieder jemand schafft, in eine Führungsposition zu kommen, der wirklich eine gute Führungspersönlichkeit ist, und nicht nur jemand, der gut reden kann oder die notwendige Portion Vitamin B hat.

Auch ist es schwierig, wenn auch nicht unmöglich, neue Parteien zu gründen und sie zu einer solchen Größe anwachsen zu lassen, dass sie echte wählbare Alternativen zu den bisher existierenden sind.

Weil nun die Demokratie unblutige Umwälzungen zulässt, sie aber auch nicht leicht macht, kann es passieren, dass zeitweilig eine personelle und inhaltliche Führungsschwäche entsteht.
Diese wird dann leicht von Populisten ausgenutzt, Menschen, die vorgeben einfache Lösungen zu haben, es in Wirklichkeit aber noch schlimmer machen.

Die Populisten können in zwei Ausführungen kommen:
\begin{enumerate}[label=\alph*)]
    \item Demokratische Populisten. Menschen mit gefährlich dummen Ideen, die aber die Demokratie nicht in Frage stellen.
    \item Undemokratische Populisten. Menschen mit gefährlich dummen Ideen, die die Demokratie in eine Autokratie umbauen wollen.
\end{enumerate}

Der Weg von a) nach b) ist nicht weit.
Wenn man merkt, dass man mit seinen dummen Ideen nicht durchkommt, kann man schnell der Demokratie die Schuld geben und von a) nach b) aufrücken, wie man das am Beispiel von Donald Trumps erster und zweiter Amtszeit gut sehen kann.

Ob, und gegebenenfalls wie, man die Demokratie noch verbessern könnte, sodass positive Umwälzungen leichter möglich wären, aber nicht gleichzeitig negative Umwälzungen erleichtert würden, weiß ich nicht.
Man kann aber versuchen, die Demokratie gegen vorübergehende Phasen von \enquote{Torenherrschaft} resilient zu machen, sodass man dem Volk ein bisschen mehr Zeit verschafft, besseres Führungspersonal hervorzubringen und gegebenenfalls einen Fehlgriff zu korrigieren.

Entscheidend ist, dass die Phasen vorübergehend sind.
Die Annahme ist, dass das Volk bei genügend Verzweiflung über die bisherigen Regierenden zwar schnell eine populistische Partei ins Amt wählen kann, aber so wie Amerika unter Trumps erster Amtszeit bald merkt, dass es einen Fehler gemacht hat.\footnote{Nach meiner Einschätzung stellt die spätere Wiederwahl von Trump keine Korrektur der Korrektur dar, sondern ist eher durch eine Kombination von Vergesslichkeit und Verzweiflung über mangelnde Wahlalternativen zustande gekommen.}

Die demokratischen Strukturen müssen dann stabil genug sein, dass eine Legislaturperiode nicht reicht, um sie umzubauen, sodass es am Ende noch freie Wahlen gibt, bei denen die Partei wieder abgewählt werden kann.
Auch wäre es gut, wenn die allerschlimmsten Entscheidungen eines populistisch beherrschten Bundestags und eines populistischen Kanzlers noch von außerhalb des Bundestages blockiert werden könnten.

Wenn die Populistenherrschaft sich über mehrere Legislaturperioden erstreckt, können sie es wahrscheinlich doch irgendwann schaffen, die Demokratie mit der Zeit nach ihren Vorstellungen umzubauen.
Es liegt in der Natur der Demokratie, dass das Volk bekommt, was es will.
Wenn es nun keine Demokratie mehr will, bekommt es die auch.

Man kann das Volk in einer Demokratie auf Dauer nicht vor sich selbst schützen.
Das geht nur in Autokratien.
Man kann es aber vor den Folgen eines Versehens bewahren.


\section{Neue Ämter und Verfassungsorgane}

\subsection{Verfassungskammer} \label{subsec:verfassungskammer}

Die Verfassungskammer wird als Ausgleich für die Proporzabweichungen durch das mehrheitssichernde Verhältniswahlverfahren gebraucht.
Sie ist der Bundestag in der Zusammensetzung nach dem ersten Wahlgang und die einzige Institution, die die Verfassung ändern kann.
\footnote{Man kann diskutieren, ob der Bundesrat weiterhin allen Verfassungsänderungen zustimmen muss. Ich gehe erst einmal davon aus, dass es so bleibt.}
Wenn die Wahl nach nur einem Wahlgang entschieden wurde, ist die Verfassungskammer mit dem Bundestag identisch.
Andernfalls tritt sie in der Zusammensetzung zusammen, die der Bundestag hätte, wenn es nur einen Wahlgang gegeben hätte.
Ihr gehören dann auch Kandidaten an, die nach dem ersten Wahlgang gewählt worden wären, aber nicht im Bundestag sitzen.
Sie kommt nur für ihre Entscheidungen zusammen.

Details siehe \mvref{subsec:verfassungsaenderungen}.

\subsection{Bundespräsidialrat} \label{subsec:bundespraesidialrat}

Es liegt in der Natur der Sache, dass vom Bürger gewählte Verfassungsorgane besonders anfällig gegen Populismus sind.
Auch nicht vom Bürger gewählte Verfassungsorgane können mit der Zeit unterwandert werden, indem sie nach und nach mit Demokratiefeinden besetzt werden.
Solche Prozesse brauchen aber Zeit.
Deswegen können nicht-gewählte Verfassungsorgane eine Zeit lang dem Populismus Grenzen setzen.

Der Gedanke ist, den Bürger die gestaltenden Verfassungsorgane (Exekutive und Legislative) selbst wählen zu lassen,
deren Arbeit aber von nicht-Bürger-gewählten Verfassungsorganen kontrollieren zu lassen.
Die kontrollierenden Verfassungsorgane dürfen selbst nicht gestalten, wie das Land sich entwickelt.
Sie sind ja auch nicht gewählt worden.
Sie sollen aber kontrollieren, dass die gewählten Gestalter ihre Arbeit richtig machen.

Ein kontrollierendes Verfassungsorgan haben wir schon.
Das Bundesverfassungsgericht führt Normenkontrollen durch.
Ihm fehlen aber Durchsetzungsmittel, und es fällt keine staatspolitischen Entscheidungen, sondern stellt nur fest, was recht ist.

Ich schlage vor, das Amt des Bundespräsidenten personell zu verstärken und mit mehr Autorität zu versehen.
Der Bundespräsident sollte einem \enquote{Bundespräsidialrat} \footnote{\enquote{Staatsrat} wäre auch ein passender Begriff. Der Bundespräsident als Staatsoberhaupt steht dem Staatsrat vor.
Der Begriff ist nur schon vorbelastet.} vorstehen,
bestehend aus ihm selbst und sechs anderen \enquote{Bundespräsidialräten}.

Der Bundespräsident leitet die Sitzungen und bestimmt was auf die Tagesordnung kommt, fällt aber ansonsten gleichberechtigt mit den anderen Ratsmitgliedern Entscheidungen.

Die personelle Verstärkung wird gebraucht, um der erhöhten staatspolitischen Verantwortung gerecht zu werden, die diesem Amt zugedacht wird.
Der Bundespräsidialrat soll in mehr Situationen als bisher der Bundespräsident allein den Bundestag auflösen oder den Kanzler absetzen können.
Außerdem sollte er Vetorechte für Exekutiv- und Legislativentscheidungen erhalten.

Der Bundespräsident behält seine repräsentativen Aufgaben.
Die Staatsoberhauptsfunktion soll aber auf sieben Räte inklusive ihm selbst aufgeteilt werden.

Der Bundespräsidialrat ist sozusagen eine \enquote{Metaregierung}.
Er regiert selbst nicht, setzt aber die Regierung im Auftrag des Volkes ein und kann sie bei Bedarf auch wieder absetzen.

\subsubsection{Besetzung des Bundespräsidialrats} \label{subsubsec:besetzung_bundespraesidialrat}

Ein kontrollierendes Verfassungsorgan ist nur so gut wie seine Unabhängigkeit von denen, die es kontrollieren soll.
Wie man parteipolitisch neutrale Organe so besetzt, dass eine Mehrheit sie weder schnell mit eigenen Leuten füllen noch durch Blockade lahmlegen kann, ist ein bekanntes Problem, für das es bewährte Lösungen gibt.\footnote{Zur Theorie solcher Organe z.~B. Bruce Ackerman, \enquote{The New Separation of Powers}, Harvard Law Review 113 (2000), S.~633 ff.; Tarunabh Khaitan, \enquote{Guarantor Institutions}, Asian Journal of Comparative Law 16 (2021), S.~S40 ff.; Mark Tushnet, \enquote{The New Fourth Branch}, Cambridge University Press 2021.
Die Empfehlungen der Venedig-Kommission des Europarats zu qualifizierten Mehrheiten und zur Auflösung von Wahlblockaden sind in CDL-PI(2025)023 zusammengefasst, \url{https://www.coe.int/en/web/venice-commission/-/cdl-pi-2025-023-e}.}
Als Vorbild dient die Besetzung des Bundesverfassungsgerichts nach Art.~93 GG und §§~6, 7 und 7a BVerfGG, die 2024 gerade zum Schutz gegen Populisten teilweise ins Grundgesetz übernommen wurde:
Wahl durch zwei Organe mit Zweidrittelmehrheit, zwölf Jahre Amtszeit ohne Wiederwahl, keine Zugehörigkeit zu Parlamenten oder Regierungen, Weiterführung im Amt bis zur Wahl eines Nachfolgers und Ersatzwahl durch das andere Wahlorgan, wenn die Wahl blockiert wird.
Diese Regeln übernehme ich.
Zusätzlich dürfen die Räte keine Parteiämter bekleiden und nach ihrer Amtszeit für eine Karenzzeit keine politischen Ämter übernehmen.
Im Folgenden beschreibe ich nur, was beim Bundespräsidialrat anders ist.

\runinhead{Wahlorgane}

Der Bundespräsident wird wie bisher von der Bundesversammlung gewählt.
Die sechs übrigen Räte werden je zur Hälfte von der Verfassungskammer und vom Bundesrat gewählt.
Die Verfassungskammer tritt an die Stelle des Bundestags, weil eine Regierungsmehrheit dort durch die Mehrheitssicherung leichter eine Zweidrittelmehrheit erreichen könnte.
Die Räte dürfen auch der Verfassungskammer nicht angehören.

\runinhead{Staffelung}

Alle zwei Jahre wird genau ein Sitz neu besetzt, abwechselnd durch Verfassungskammer und Bundesrat.
Bei der Ersteinrichtung werden die Amtszeiten dafür einmalig auf zwei, vier, sechs, acht, zehn und zwölf Jahre festgelegt.
Eine Mehrheit kann so in einer Legislaturperiode höchstens zwei der sechs Sitze neu besetzen.
Für eine eigene Mehrheit von vier Stimmen bräuchte sie acht Jahre, mit einem eigenen Bundespräsidenten sechs Jahre, und das jeweils mit Zweidrittelmehrheit in Verfassungskammer und Bundesrat.

\runinhead{Tagesordnung}

Das Recht des Bundespräsidenten, die Tagesordnung zu bestimmen, darf nicht dazu führen, dass er unliebsame Themen blockieren kann.
Zwei Räte gemeinsam können deshalb jederzeit eine Sitzung verlangen und einen Punkt auf die Tagesordnung setzen lassen, der dann zeitnah behandelt werden muss.

\runinhead{Ausfall von Räten}

Die Weiterführung im Amt hilft gegen eine Blockade nur, solange die bisherigen Räte ihr Amt noch ausüben können.
Stirbt ein Rat oder tritt er zurück, könnte eine Sperrminorität den Rat nach und nach ausdünnen, bis er handlungsunfähig ist.
Die Ersatzwahl durch das andere Wahlorgan hilft nur, wenn die Blockade auf eines der beiden beschränkt ist.
Außerdem würde jedes vorzeitige Ausscheiden die Staffelung durchbrechen und der gerade amtierenden Mehrheit einen zusätzlichen Sitz verschaffen.

Ich schlage deshalb vor, jeden Rat zusammen mit einem persönlichen Stellvertreter zu wählen, in derselben Wahl und mit derselben Mehrheit.
Fällt der Rat aus, rückt der Stellvertreter für den Rest der Amtszeit nach.
Der Ersatz wurde damit gewählt, als noch Einigkeit bestand, und eine spätere Blockade kann ihn nicht mehr verhindern.

Der Rat bleibt beschlussfähig, solange mindestens fünf Räte im Amt sind.
Eine Entscheidung braucht aber immer vier Stimmen, auch wenn weniger Räte im Amt sind.
Als letzte Notlösung wählen die verbleibenden Räte einen frei gebliebenen Sitz bis zur nächsten erfolgreichen regulären Wahl selbst nach.
Das muss die Ausnahme bleiben, weil sich der Rat sonst selbst ergänzen und von jeder demokratischen Bindung lösen würde.

\runinhead{Abberufung}

Ein Bundespräsidialrat kann nur durch das Bundesverfassungsgericht abberufen werden.
Den Antrag dazu können Verfassungskammer oder Bundesrat jeweils mit Zweidrittelmehrheit stellen, ähnlich wie bei der Präsidentenanklage nach Art.~61 GG.
Gründe für eine Abberufung sind eine vorsätzliche Verletzung des Grundgesetzes oder eines Bundesgesetzes, ein schweres Vergehen im Sinne von Abschnitt~\ref{einzelne_parlamentarier_verstossen_gegen_gesetz}, die dauerhafte Unfähigkeit, das Amt auszuüben, oder wenn der Rat nach seinen Zielen oder seinem Verhalten erkennbar darauf ausgerichtet ist, die Demokratie abzuschaffen oder in eine Ordnung umzubauen, die dem freiheitlichen Grundgedanken des Grundgesetzes zuwiderläuft.

\runinhead{Verankerung im Grundgesetz}

Beim Bundesverfassungsgericht steht die Zweidrittelmehrheit nur im Bundesverfassungsgerichtsgesetz und könnte mit einfacher Mehrheit abgeschafft werden.
Beim Bundespräsidialrat sollten deshalb alle Besetzungsregeln im Grundgesetz stehen, auch die Zweidrittelmehrheit sowie die Stellvertreter- und Ausfallregelung.
Nur Fristen und Verfahrensfragen können einem Bundesgesetz überlassen werden.

\runinhead{Gegenseitige Absicherung und Restrisiko}

Auch wer den Rat beherrscht, kann nicht nach Belieben Bundeskanzler und Bundestag absetzen.
Er braucht dafür entweder den gemessenen Volkswillen (\ref{subsubsubsec:vorgezogene_neuwahlen_auf_volkswunsch_hin}) oder eine stichhaltige Begründung, die das Bundesverfassungsgericht überprüfen kann (\ref{absetzung_durch_bundespraesidialrat}).
Ein Missbrauch setzt also voraus, dass auch das Bundesverfassungsgericht unterwandert ist.
Das unterscheidet den Bundespräsidialrat von nicht gewählten Organen in sogenannten Vormundschaftsdemokratien (\emph{tutelary democracy}) wie Thailand, die gewählte Regierungen nach politischem Gutdünken absetzen.

Schwerer zu kontrollieren ist ein unterwanderter Rat, der einfach nichts tut, also kein Veto einlegt und keine Neuwahlen anordnet.
Gegen eine Mehrheit, die über viele Jahre Zweidrittelmehrheiten in Verfassungskammer und Bundesrat erringt, hilft letztlich keine Besetzungsregel, wie Ungarn und Polen gezeigt haben.
Die Staffelung verschafft dem Volk aber die Zeit, die nach Abschnitt~\ref{sec:haertung_der_demokratie} gebraucht wird, um einen Fehlgriff bei der nächsten Wahl zu korrigieren.

\subsection{Bundesamt für Meinungsforschung}\label{subsec:bundesamt_fuer_meinungsforschung}

Für alle Regierungen ist es wichtig zu wissen, was das Volk denkt.
Demokratische Regierungen bilden da keine Ausnahme.
Für den in \ref{subsec:demokratische_notbremse} gemachten Vorschlag, den Volkswillen bei der Entscheidung für vorgezogene Neuwahlen mitzuberücksichtigen,
ist ein zuverlässiges Stimmungsbild der Bevölkerung essentiell.
Man kann sich bei solchen Entscheidungen nicht allein auf Umfragen privater Meinungsforschungsinstitute verlassen.

Ich schlage deshalb vor, ein Bundesamt für Meinungsforschung einzurichten, das ständig damit beschäftigt ist, mit nicht-geheimdienstlichen Methoden,
ein anonymisiertes Meinungsbild der Bevölkerung zu bestimmten Fragen zu erstellen und aktuell zu halten.
Es sollte sich dabei auf eine größere Datenbasis stützen als die kostengünstigen repräsentativen Umfragen privater Meinungsforschungsinstitute.
Deren Ergebnisse sollten aber auf jeden Fall zur Gegenkontrolle mitherangezogen werden.

\runinhead{Unabhängigkeit}
Das Bundesamt für Meinungsforschung ist fachlich unabhängig, ähnlich wie das Statistische Bundesamt.
Das Amt legt seine Methodik offen.

\runinhead{Auftraggeber}
Umfragen in Auftrag geben können die Bundesregierung, die Fraktionen des Bundestags und der Bundespräsidialrat.
Regierung und Opposition müssen dabei gleichermaßen Zugang zum Amt haben, nicht nur die Regierung allein.
Da auch ein großes Amt nur eine begrenzte Kapazität hat, braucht es dafür Kontingente.
Die regelmäßige Erhebung des Neuwahlwunsches und die Umfragen des Bundespräsidialrats dürfen davon nicht abhängen.
Der Auftraggeber bestimmt das Thema, das Amt den Wortlaut der Fragen.
Es kann Suggestivfragen ablehnen.

\runinhead{Veröffentlichung}
Der Auftraggeber entscheidet, ob und wann die Ergebnisse seiner Umfrage veröffentlicht werden.
Dafür gibt es zwei Gründe:
\begin{enumerate}[label=\alph*)]
    \item Regierung und Opposition wollen oft wissen, wie das Volk über eine Frage denkt, bevor sie sich positionieren oder ihre Entscheidung kommunizieren.
    \item Veröffentlichte Umfragen können einen Verstärkereffekt auslösen, nach dem Motto: \enquote{Wenn X\% der Bevölkerung das so sehen, dann kann es ja nicht falsch sein.}
    Handelt es sich dabei um populistisches Gedankengut, dem man entgegenwirken möchte, macht man sich damit das Leben schwerer.
    Zahlen eines staatlichen Amtes haben dabei mehr Gewicht als die eines privaten Instituts.
\end{enumerate}

Damit sich der Auftraggeber nicht die passenden Zahlen herauspicken kann, veröffentlicht das Amt selbst, und zwar vollständig mit dem Wortlaut der Fragen, allen Antworten und der Methodik.

Immer veröffentlicht werden der regelmäßig erhobene Neuwahlwunsch und alle Ergebnisse, auf die der Bundespräsidialrat seine Entscheidungen stützt.
Eine Geheimhaltung würde hier ohnehin nicht halten, da zu viele Menschen die Zahlen kennen.
Schon das Schweigen wäre ein Signal, und es würde Verschwörungstheorien Nahrung geben.

Spätestens nach Ende der folgenden Legislaturperiode werden alle Ergebnisse veröffentlicht.
Strategisch spielen sie dann keine Rolle mehr, aber man kann nachprüfen, ob das Amt missbraucht wurde.


\section{Nachfolgeregelungen bei vorzeitiger Ablösung der Regierung oder Teilen davon}\label{sec:abloseregelungen}

Aus vielen Gründen kann es vorkommen, dass Bundestag und Bundesregierung oder auch nur Teile davon
vorzeitig abgelöst werden müssen.
Damit das Land zu keinem Zeitpunkt ohne Regierung bleibt, muss es für alle denkbaren Fälle Ablöseregelungen geben.
Die aktuell gültigen Regeln können nicht unmodifiziert übernommen werden.
Dazu sind die vorgeschlagenen Verfahrensänderungen zu grundlegend.
Zum Beispiel soll der Bundeskanzler jetzt indirekt vom Volk gewählt werden und nicht mehr vom Bundestag.
Daher kann der Bundestag nicht einfach so einen neuen Bundeskanzler wählen.
Das käme nur als befristete Übergangslösung in Betracht.

Im Folgenden werden daher die denkbaren Fälle einzeln durchgegangen:

\subsection{Ablösung des Bundeskanzlers ohne Legislativänderung}\label{subsec:ablosung-des-bundeskanzlers-ohne-legislativanderung}
Es kann vorkommen, dass der Bundeskanzler vorzeitig abgelöst werden muss, obwohl sich an den Mehrheitsverhältnissen im Parlament
ansonsten nichts geändert hat.

Der Bundeskanzler könnte z. B. schwer erkranken oder sogar sterben, er könnte von sich aus zurücktreten
oder sein Amt nicht mehr ausüben dürfen, weil er sich persönlich etwas so Schwerwiegendes hat zuschulden kommen lassen,
dass er nicht mehr Bundeskanzler sein darf.
In diesem Fall müssen keine Parlamentsneuwahlen durchgeführt werden.
Es reicht, den Kanzler neu zu wählen.

Die Partei, der der abzulösende Kanzler entstammt, muss wieder drei Kanzlerkandidaten dem Volk zur Wahl
vorstellen.
Zwei der drei Kandidaten dürfen aber müssen nicht dieselben sein, wie bei der ursprünglichen Bundestagswahl.
Bei einer Einparteienregierung wird der Kandidat mit den meisten Stimmen automatisch Bundeskanzler.
Der Bundespräsident muss ihn ernennen.

Bei einer Koalitionsregierung greift das Prozedere aus dem \mvref{sec:wahl-der-spitzenkandidaten-durch-das-volk}.
Der Bundeskanzler kann nach Methode a) oder b) bestimmt werden und bedarf einer Bestätigung durch die Koalitionäre
in geheimer Wahl.

Erfolgt diese nicht, kommt noch eine Koalitionsneubildung analog zu \ref {koalitionszerbruch-mit-neubildung} infrage.
Der einzige Unterschied zu \ref {koalitionszerbruch-mit-neubildung} ist, dass Option c) naturgemäß nicht mehr zur
Verfügung steht.
Sollte die alte Kanzlerpartei Teil der neuen Koalition sein, geht sie mit dem Ergebnis der Kanzlerkandidatennachwahl \footnote{Der Begriff \enquote{Kanzlerneuwahl} wäre nicht korrekt, da es ja auch von den Wahlergebnissen der anderen Koalitionspartner abhängt, ob überhaupt einer von den neuen Kanzlerkandidaten Bundeskanzler wird.} in diese Koalition.

Ist eine Koalitionsneubildung nicht gewünscht oder auch nicht erfolgreich, müssen Bundestagsneuwahlen stattfinden.
Bis zur Ernennung eines neuen Bundeskanzlers durch den Bundespräsidenten führt der Vizekanzler kommissarisch die Regierungsgeschäfte.

\runinhead{Hinweis:}
Auch wenn von vornherein klar ist, dass die bestehende Regierungskoalition an dem fehlenden Kanzler zerbrechen wird,
muss trotzdem eine Kanzlerkandidatennachwahl stattfinden, sofern die alte Kanzlerpartei an einer neuen Koalition teilnimmt.
Das gilt selbst, wenn die Teilnahme an einer neuen Koalition nicht unmittelbar nach dem Auseinanderbrechen der alten erfolgen sollte,
sondern zum Beispiel erst nach dem Auseinanderbrechen gegebenenfalls mehrerer neuer Koalitionen.
Der Grund ist, dass sie keinen vom Volk aus drei Kandidaten ausgewählten Kanzlerkandidaten mehr vorweisen kann.

Der Bundespräsident moderiert diesen Prozess.
Er ordnet eine Kanzlerkandidatennachwahl an, oder verzichtet darauf, wenn sich eine Koalition ohne Beteiligung der alten Kanzlerpartei findet,
oder ordnet neue Parlamentswahlen an, sobald klar ist, dass sich keine Nachfolgekoalition finden wird.
Im Falle einer Einparteienregierung, die eben nicht auf eine Koalition zum Regieren angewiesen ist, ordnet der Bundespräsident
im Falle eines reinen Kanzlerausfalls immer eine Kanzlerkandidatennachwahl an.

Auch eine Einparteienregierung kann als Ganzes abgelöst werden müssen, weil sich z.B. die Partei zerstreitet.
Dieser Fall ist wie der Zerbruch einer Koalition zu behandeln und von der hier beschriebenen Ablösung der Person des Bundeskanzlers zu unterscheiden.

\subsection{Zerbruch einer Koalition mit anschließender Bildung einer neuen Koalition} \label{koalitionszerbruch-mit-neubildung}
Es passiert ja häufig, dass Koalitionen zerbrechen und manchmal finden sich hinterher auch neue zusammen.
Wenn sich eine neue Koalition zusammengefunden hat, muss sie dies dem Bundespräsidenten mitteilen, damit er Kanzlerkandidatennachwahlen anordnet,
in dem Fall, dass eine Partei keine drei vom Volk gewählte Kanzlerkandidaten vorweisen kann.
Siehe \ref {subsec:ablosung-des-bundeskanzlers-ohne-legislativanderung}.

Nachdem alle Koalitionsparteien mit vom Volk gewählten Kanzlerkandidaten in die Koalition gehen, ist der Test für eine
erfolgreiche Koalitionsbildung die Kanzlerwahl durch die Koalition gemäß dem \mvref{sec:wahl-der-spitzenkandidaten-durch-das-volk}.
Im Unterschied zu Koalitionsbildungen direkt nach einer Bundestagswahl sollte es aber zusätzlich die Möglichkeit geben, dass
der alte Bundeskanzler im Amt bleibt, wenn die Koalitionäre sich darauf einigen und der alte Bundeskanzler für eine Wahl
zur Verfügung steht.

Das gilt selbst, wenn der alte Bundeskanzler keiner Partei der neuen Koalition angehören sollte, also ohne eigene Partei
im Rücken regieren würde.
Der Sinn dieser Regelung ist maximale Stabilität der Exekutive zu ermöglichen.
Wenn der Glücksfall eintritt, dass man einen Bundeskanzler hat, der in seinem Amt erfahren ist und als Person von allen
Parteien respektiert wird, dann soll er sein Amt trotz Parteienstreitigkeiten weiter ausüben dürfen, wenn er selbst das
will und die neuen Koalitionäre dies auch wollen.

Die Koalitionäre haben also die Wahl zwischen Kandidaten, die nach drei verschiedenen Methoden ermittelt wurden:
\begin{enumerate}[label=\alph*)]
  \item Derjenige Kandidat der stärksten Koalitionspartei, der die meisten Stimmen erhalten hat.
  \item Derjenige Kandidat von allen Koalitionsparteien, der die meisten Stimmen erhalten hat.
  \item Der alte Bundeskanzler.
\end{enumerate}

Nach diesen drei Methoden können drei verschiedene Kandidaten herauskommen, oder auch nur zwei oder einer, weil z. B
die alte Kanzlerpartei auch die stärkste Koalitionspartei in der neuen Koalition ist, und der alte Bundeskanzler auch
der Kandidat, der bei der letzten Bundestagswahl absolut die meisten Stimmen auf seine Person vereinen konnte.

Sollte bei der geheimen Kanzlerwahl (ggf. in mehreren Wahlgängen) kein Bundeskanzler bestätigt werden, ist die Koalition
nicht zustande gekommen und der Bundespräsident ordnet Parlamentsneuwahlen an, sobald klar ist, dass keine andere Koalition
zustande kommen wird und der bisherige Bundeskanzler auch keine Minderheitsregierung bilden wird.

\subsection{Zerbruch einer Koalition ohne Bildung einer neuen Koalition}
Wenn eine den Bundeskanzler tragende Koalition zerbricht, und keine neue gebildet wird, heißt dies nicht unbedingt, dass
der Bundeskanzler deswegen zurücktreten müsste.
Er kann sein exekutives Amt auch ohne eigene Mehrheit in der Legislative weiter ausüben, solange er nicht explizit abgesetzt wird.
Wenn er die Zustimmung des Parlaments für seine Vorhaben braucht, muss er sich eben je nach Vorhaben passende Mehrheiten suchen.

Wenn eine Regierungskoalition zerbricht, sich keine neue dafür findet und der alte Bundeskanzler auch keine Minderheitsregierung
führen will, muss der Bundespräsident Parlamentsneuwahlen anordnen.
Der alte Bundeskanzler muss dies dem Bundespräsidenten mitteilen.

\subsection{Ablösung einzelner Minister}
Das Verfahren zur Ablösung der Bundesminister bleibt wie gehabt: Der Bundespräsident ernennt oder entlässt sie auf Vorschlag des Bundeskanzlers.

\subsection{Ablösung einzelner Parlamentarier} \label {abloesung_einzelner_parlamentarier}
Wenn einzelne Parlamentarier abgelöst werden müssen, rücken bei Listenkandidaten einfach die nächsten nach.
Bei Direktwahlkandidaten muss eine Personennachwahl im zugewiesenen Wahlkreis durchgeführt werden.


\section{Demokratisches Gegengewicht} \label{sec:demokratisches-gegengewicht}
Durch das vorgeschlagene Wahlverfahren bekommt die Regierung mehr Macht.
Dafür muss an anderer Stelle ein demokratischer Ausgleich geschaffen werden, damit die Gefahr des Festsetzens an der Macht wieder verringert wird.

Die Regierung hat mehr Macht bekommen, damit sie ihren Auftrag, zu gestalten, auch in herausfordernden Zeiten, wo viel entschieden werden muss, erfüllen kann.
Umso wichtiger ist es, dass es eine \enquote{demokratische Notbremse} gibt, mit der man eine schlechte Regierung auch vorzeitig wieder loswerden kann, wenn diese ihre Macht missbraucht oder durch Inkompetenz verschwendet.
Es muss Mechanismen geben, eine vorzeitige Ablösung von Einzelpersonen oder der ganzen Regierung einzuleiten, wenn dies geboten ist.

Die Herausforderung ist, das Machtgleichgewicht so zu gestalten, dass eine Regierung einerseits genug Macht hat, um auch unpopuläre Entscheidungen durchzusetzen, und andererseits noch gestoppt werden kann, wenn sie zu viel Schaden anrichtet.
Auch in einer Demokratie muss es möglich sein, zeitweilig gegen den Willen der Mehrheit der Bevölkerung zu regieren.
Wenn zu jeder Zeit der Volkswille geschieht, hat man eigentlich gar keine Regierung.
Manche Maßnahmen sind einfach notwendig, obwohl sie keiner mag.

Der Unterschied zu einer Diktatur ist, dass Regieren gegen das Volk nur zeitweilig möglich ist und persönlichen Einsatz erfordert.
Eine notwendige, aber unpopuläre Maßnahme kann einen die Wiederwahl kosten.
Das soll auch so sein, weil es die Wahrscheinlichkeit von Machtmissbrauch verringert.

Das Problem ist aber, dass manche Maßnahmen nicht nur unpopulär sind, sondern so gravierende Fehler gemacht werden, dass man nicht bis zum Ende der Legislaturperiode warten kann, um sie zu korrigieren.

In den folgenden Abschnitten werden daher verschiedene Ablösungsszenarien diskutiert.
Während es in Kapitel~\ref{sec:abloseregelungen} um Nachfolgeregelungen ging (Was passiert, wenn aus irgendeinem Grund die Regierung oder Teile davon vorzeitig abgelöst werden müssen?),
geht es in diesem Abschnitt darum, wie und unter welchen Umständen eine vorzeitige Ablösung eingeleitet werden kann.

\subsection{Begriffsklärung: Mehrheitsstufen} \label{mehrheitsstufen}
Vor der eigentlichen Diskussion zunächst eine Begriffsklärung.
Man kann verschieden stark qualifizierte Mehrheiten definieren, z.~B.:

\begin{table}[h]
\renewcommand{\arraystretch}{1.4}%
\begin{tabularx}{\linewidth}{| c | l | X | p{3.5cm} |}
  \hline
  \textbf{Stufe} & \textbf{Name} & \textbf{Bedeutung} & \textbf{Prozentsätze} \\
  \hline
  1 & Einfache Mehrheit    & Mehr Menschen sind dafür als dagegen. & $> 50\%$ \\
  \hline
  2 & Deutliche Mehrheit   & Es ist klar erkennbar, dass mehr Menschen dafür als dagegen sind. Zufällige Effekte können ausgeschlossen werden. & ${\approx}\,60\%$ (kontextabhängig) \\
  \hline
  3 & Zweidrittel Mehrheit & Doppelt so viele Menschen sind dafür, wie dagegen. & $66{,}66\%$ \\
  \hline
  4 & Übergroße Mehrheit   & Viel mehr Menschen sind dafür als dagegen. & $80\%\ldots90\%$ \\
  \hline
  5 & Fast alle            & Fast alle Abstimmenden sind dafür. Nur einzelne Ausnahmen dagegen. & $99{,}9\%\ldots99{,}999\%$ (kontextabhängig) \\
  \hline
  6 & Einstimmig           & Alle Abstimmenden sind dafür. & $100\%$ \\
  \hline
\end{tabularx}
\caption{Mehrheitsstufen}
\label{tab:mehrheitsstufen}
\end{table}

Wenn ich mir irgendwo nicht sicher bin, wie stark die erforderliche Mehrheit qualifiziert sein sollte, dann schreibe ich einfach \enquote{n-te Mehrheit} oder \enquote{n-stufige Mehrheit}.
Damit ist eine Mehrheit n-ter Stufe gemäß obiger Tabelle gemeint.
D.~h., die genaue Qualifizierung muss noch festgelegt werden.

Die \enquote{Mehrheitsstufe} ist ein Parameter, den man hat, um das im vorigen Abschnitt angesprochene Verhältnis zwischen notwendiger Autorität der Regierung und ihrer Korrigierbarkeit auszubalancieren.
Die Einteilung der Mehrheitsstufen gemäß Tabelle~\ref{tab:mehrheitsstufen} ist mehr als ein Beispiel zu verstehen, um das Prinzip zu verdeutlichen, als eine endgültige Definition.
Man kann auch noch mehr Stufen definieren oder welche weglassen, wenn sich zeigt, dass man sie sowieso nicht braucht.

Manchmal werde ich nur von \enquote{hohen} und \enquote{niedrigen} Mehrheiten sprechen, ohne festzulegen, welche Mehrheitsstufen in diesem Zusammenhang genau damit gemeint sind.
Die Festlegung muss dann in dem Diskussionsprozess erfolgen, den ich hoffe, mit diesem Dokument anzustoßen.

\subsection[]{Einzelne Parlamentarier verstoßen gegen das Gesetz} \label{einzelne_parlamentarier_verstossen_gegen_gesetz}

Einzelne Parlamentarier können bereits jetzt nach Art 46 GG strafrechtlich verfolgt werden, wenn der Bundestag dem zustimmt.
Darüber hinaus sollte der Bundestag das Recht haben, Abgeordnete für die Dauer des Verfahrens zu suspendieren und bei
nachgewiesenen schweren Vergehen das Mandat sogar rückwirkend zu entziehen, sodass der Abgeordnete alle Privilegien
verliert, die mit einem Bundestagsmandat verbunden sind.

Schwere Vergehen in diesem Sinne sind Vergehen
\begin{itemize}
  \item die sich vorsätzlich gegen die Bundesrepublik Deutschland richten, wie z.~B. Spionage für eine fremde Macht, oder
  \item Vergehen, die nach allgemeinem Rechtsempfinden ein Volksvertreteramt ausschließen, wie z.~B. Gewaltverbrechen oder Sexualstraftaten.
 \end{itemize}

Welche Vergehen genau nach allgemeinem Rechtsempfinden ein Volksvertreteramt ausschließen, muss noch bestimmt werden.
Der Grundsatz sollte jedoch klar sein.

99\% der Bevölkerung würden sich wohl nicht von einem Mörder oder Kinderschänder vertreten lassen wollen, auch wenn er anderweitig noch so brillant ist.
Bei Steuerhinterziehung und Trunkenheitsfahrten hingegen könnten die Meinungen vielleicht geteilt sein.

Der Bundestag hat das Recht bei Nachweis eines schweren Vergehens, das Mandat ab sofort zu entziehen oder sogar rückwirkend, wenn der Strafbestand schon die ganze
Legislaturperiode bestanden hat.
Der betroffene Abgeordnete hat das Recht, bei dem Bundesverfassungsgericht gegen diese Maßnahme zu klagen.
Gibt es umgekehrt im Bundestag keine Mehrheit für die Suspendierung oder den Mandatsentzug, können einzelne Abgeordnete beim Bundesverfassungsgericht darauf klagen, und das
Bundesverfassungsgericht kann der Klage entsprechen oder nicht.

\subsection{Absetzung des Bundeskanzlers}

\subsubsection{Denkbare Gründe für den Wunsch nach einer vorzeitigen Absetzung des Bundeskanzlers}

Es sind verschiedene Gründe für den Wunsch nach Absetzung des Bundeskanzlers denkbar.
Nicht alle davon sollten tatsächlich zu einer Absetzung führen.

\begin{enumerate}
    \item Inkompetenz. Der Bundeskanzler kann sich einfach als unfähig erweisen.
    \item Der Bundeskanzler trifft Entscheidungen, die von seiner Partei/Fraktion nicht mehr mitgetragen werden können.
    \item Eine Koalition zerbricht und es bildet sich eine neue dafür.
    \item Der Bundeskanzler trifft Entscheidungen, die von der n-ten Mehrheit des Volkes nicht gut gefunden werden.
    \item Der Bundeskanzler trifft Entscheidungen, die das Land mittel- bis langfristig ins Verderben führen.
    \item Der Bundeskanzler trifft Entscheidungen, die das Land kurzfristig in eine Katastrophe stürzen. Man muss ihn sofort aufhalten.
    \item Der Bundeskanzler verstößt gegen das Gesetz.
        Zu unterscheiden sind
        \begin{enumerate}
          \item Verstöße, die aufgrund ihrer Schwere eine Fortsetzung der Kanzlerschaft unmöglich machen. Hier gelten dieselben Erwägungen wie bei allen Parlamentariern nur etwas strenger, da der Kanzler ja viel mehr Macht und auch eine größere Vorbildfunktion hat, als einfache Parlamentarier.
            \begin{enumerate}
                \item Der Kanzler hält sich nicht an die Verfassung bzw. befolgt Entscheidungen des Bundesverfassungsgerichts nicht.
                \item Der Kanzler begeht Hochverrat, indem er z.~B. vorsätzlich Staatsgeheimnisse verrät.
                \item Der Kanzler begeht ein Verbrechen wie z.~B. Mord, das ihn nach allgemeinem Rechtsempfinden für eine Kanzlerschaft disqualifiziert.
            \end{enumerate}
          \item Verstöße, die während der Kanzlerschaft geahndet werden können, ohne die Kanzlerschaft zu beeinträchtigen, wenn z.~B. nur eine Geldstrafe fällig ist.
          \item Verstöße, die nicht so schwerwiegend sind, dass die Kanzlerschaft nicht fortgesetzt werden könnte, aber erst nach der Kanzlerschaft geahndet werden können. (Z.B. weil eine Gefängnisstrafe droht, die aber nicht sofort vollstreckt werden muss.)
        \end{enumerate}

    \item Der Bundeskanzler verstößt zwar gegen kein Gesetz, aber macht etwas, das
        \begin{enumerate}[label=\alph*)]
            \item die (n-te) Mehrheit des Bundestags ihn seines Amtes entheben will
            \item die (m-te) Mehrheit des Volkes ihn seines Amtes entheben will
        \end{enumerate}
    \item Der Bundeskanzler hält seine Wahlversprechen nicht.
    \item Der Bundeskanzler macht schlechte Politik
\end{enumerate}

Die letzten beiden Gründe sind allein nicht ausreichend für eine vorzeitige Absetzung des Bundeskanzlers, sondern traurigerweise normal in einer Demokratie.
Wenn man einen besseren Bundeskanzler hat, hat man Glück.
Die Quittung für gebrochene Wahlversprechen und durchschnittlich schlechte Politik kommt bei der nächsten Wahl.
Die anderen Punkte verdienen eine Diskussion.

\subsubsection{Mittel gegen Inkompetenz: Probezeit für den Bundeskanzler}
Eine Besonderheit des Bundeskanzleramts ist, dass es eigentlich kein vergleichbares nationales Amt gibt, was man als Referenz für die Eignung einer Person zum Bundeskanzler heranziehen könnte.
Das nächstliegende Amt ist noch das Ministerpräsidentenamt eines Bundeslandes.
Aber auch das ist nicht wirklich vergleichbar.
Eine Person kann in einem Bundesland ein beliebter \enquote{Landesvater} sein, aber trotzdem völlig ungeeignet, das ganze Land zu führen.

Die einzig wahre Möglichkeit, die Kanzlerfähigkeiten einer Person zu testen, ist, sie das Amt eine Zeit lang ausüben zu lassen.
Ich schlage deshalb vor, eine Probezeit für einen neuen Bundeskanzler einzuführen.
Wenn jemand als Bundeskanzler noch keine Probezeit bestanden hat, erhält er das Amt erst einmal auf Probe.
Das bedeutet, dass er nach einer noch festzusetzenden Frist (Vorschlag: halbe Legislaturperiode), vom Volk im Amt bestätigt werden muss, damit er Bundeskanzler bleiben darf.
Wie in jedem anderen Beruf ist die Hürde für eine Trennung während der Probezeit deutlich niedriger als danach.
Anders als bei vorgezogenen Neuwahlen auf Volkes Wunsch (\ref{subsubsubsec:vorgezogene_neuwahlen_auf_volkswunsch_hin}) ist deshalb keine Zustimmung des Bundespräsidialrats erforderlich.

Es wird eine Volksabstimmung durchgeführt mit der Frage, ob der amtierende Bundeskanzler sein Amt bis zum Ende der Legislaturperiode ausüben darf.
Antwortet eine n-te Mehrheit mit \enquote{nein}, hat er die Probezeit nicht bestanden und das Verfahren zur Ablösung des Bundeskanzlers nach Abschnitt~\ref{subsec:ablosung-des-bundeskanzlers-ohne-legislativanderung} kommt zur Anwendung.
Andernfalls hat er seine Probezeit bestanden und muss nie wieder auf Probe regieren.
Das gilt auch, sollte er zwischenzeitlich abgewählt und später erneut gewählt werden.
Das Volk kennt ihn jetzt im Amt und kann eine bewusste Entscheidung treffen.

Für n würde ich 2, d.h.~eine deutliche Mehrheit, vorschlagen.
Eine einfache Mehrheit wäre zu wenig, weil auch ein neuer Bundeskanzler schon in seiner ersten Amtszeit unpopuläre Entscheidungen durchsetzen muss.
Die Mehrheitsstufe sollte aber niedriger sein als bei vorgezogenen Neuwahlen auf Volkes Wunsch, damit man einen ungeeigneten Bundeskanzler in der Probezeit vergleichsweise leicht wieder loswird.

Man kann diskutieren, ob nach einer nicht bestandenen Probezeit eine erneute Kanzlerkandidatur nach einer Sperrfrist wieder möglich sein sollte oder nicht.

\runinhead{Gründe dagegen:}

\begin{itemize}
  \item Stabilität des politischen Systems. Das Volk kennt ihn schon. Warum sollte er es beim nächsten Mal besser als bei diesem Mal machen? Wenn er beim nächsten Versuch wieder durchfällt, muss wieder neu gewählt werden und so viele vorfristige Wahlen sind organisatorischer Aufwand und tun der Stabilität des Systems nicht gut.
  \item Das Volk vergisst schnell. Es besteht die Gefahr, dass ein völlig ungeeigneter Kandidat, wie Donald Trump, erneut an die Macht kommt und beim zweiten Mal schlimmer als beim ersten Mal wird.
 \end{itemize}

\runinhead{Gründe dafür:}
\begin{itemize}
  \item Menschen können sich tatsächlich ändern. Wenn die Sperrfrist lange genug ist, sagen wir drei Legislaturperioden, was in Deutschland 12 Jahren entsprechen würde, kann der Kandidat inzwischen so gereift sein, dass er sogar ein besonders guter Bundeskanzler wird, weil er aus seinen ersten Erfahrungen gelernt hat.
  \item Kandidatenmangel. Es ist schwer, gutes Führungspersonal zu finden. Vielleicht kann man es nach ausreichend Sperrzeit noch einmal mit einem schon bekannten Kandidaten probieren.
 \end{itemize}

\subsubsection{Absetzung des Bundeskanzlers durch den Bundestag}
Wenn der Bundeskanzler sich von der Linie seiner Partei in einigen Fragen entfernt, ist dies noch lange kein Grund, ihn abzusetzen.
Er ist schließlich, wie alle anderen Abgeordneten auch, in erster Linie dem Volk und seinem Gewissen verpflichtet und nicht seiner Partei.

Wenn sie ihn dafür bei der nächsten Wahl nicht mehr als Kandidat aufstellen wollen, ist das ihr gutes Recht.
Während einer Legislaturperiode können sie ihn aber nicht von seinem Posten abberufen.

Wenn die Fraktion sich mit dem Bundeskanzler in einigen Fragen überwirft, ist das auch noch kein Grund ihn abzusetzen.
Sie müssen seinen Vorhaben ja nicht zustimmen.
Der Bundeskanzler muss dann sehen, ob er die fehlende Unterstützung vielleicht bei Teilen der Opposition bekommt.
Wenn ja, kann er seine Entscheidungen vielleicht doch noch durchsetzen.
Wenn nein, muss er eben nachgeben.
Das ist ein normaler demokratischer Prozess.

Nichtsdestotrotz kann es Fälle geben, wo die Fraktion nicht nur in einzelnen Fragen mit dem Bundeskanzler nicht übereinstimmt, sondern den Kanzler loswerden will.
Wenn sie das im Rahmen einer Koalitionsneubildung tut, ist das Verfahren bereits in Abschnitt ~\ref {koalitionszerbruch-mit-neubildung} geregelt.

Nach dem aktuell gültigen Grundgesetz könnte die Fraktion jederzeit den Bundeskanzler gegen ein anderes Mitglied des Bundestages austauschen, indem sie nach Artikel 67 dem aktuellen Bundeskanzler das Misstrauen ausspricht und einen anderen Bundeskanzler an seiner Stelle wählt.
Da nach Artikel 67 die einfache Mehrheit des Bundestags für ein Misstrauensvotum reicht, könnte die Fraktion theoretisch jederzeit z.~B. Friedrich Merz gegen Lars Klingbeil austauschen, wenn dem nur fast alle Mitglieder der Koalition zustimmen würden.

Nach dem neuen System wäre ein Austausch von Personen nicht möglich, außer er käme als Nebenprodukt einer Koalitionsneubildung zustande.
Allenfalls Kanzlerkandidatennachwahlen könnten ausgelöst werden.
Artikel 67 kann in dieser Form also nicht beibehalten werden.

Gerade in den gegenwärtigen, unruhigen Zeiten, wo man sich auf nichts mehr verlassen kann, muss man aber davon ausgehen, dass es Situationen geben kann, in denen ein Bundeskanzler gestoppt werden muss, weil er zu viel Schaden anrichtet.

\runinhead{Dabei sind zwei Fälle zu unterscheiden:}
\begin{enumerate}
    \item Der Bundeskanzler muss mit sofortiger Wirkung seines Amtes enthoben werden.
    \item Der Bundeskanzler kann noch geschäftsführend im Amt bleiben, bis der Nachfolger gewählt wurde.
\end{enumerate}

Im ersten Fall kann der Bundestag mit einer n-stufigen Mehrheit den Bundeskanzler seines Amtes entheben und auch gleich einen Interimsnachfolger einsetzen, wenn er zu dem Schluss kommen sollte, dass der Vizekanzler die Amtsgeschäfte auch nicht weiterführen kann, weil er z.~B. in denselben Skandal verwickelt ist.
Auf jeden Fall werden Kanzlerkandidatennachwahlen ausgelöst, die innerhalb einer noch festzusetzenden Frist stattfinden müssen.

Im zweiten Fall reicht eine m-stufige Mehrheit aus, um den Bundeskanzler seines Amtes zu entheben und die Kanzlerkandidatennachwahlen auszulösen.

Für n würde ich 3, d.~h. eine zweidrittel Mehrheit vorschlagen, da es sich doch um einen sehr schwerwiegenden Eingriff in die Staatsführung handelt.
Für m würde ich 2, d.~h. eine deutliche Mehrheit wählen.
M=1, wie nach der aktuellen Regelung wäre aber auch möglich.
Dann könnte die Fraktion, in Eigenregie den Bundeskanzler stürzen, wenn sie sich denn einig wäre.
Nur einen Nachfolger könnte sie nicht selbst einsetzen.

\subsubsection{Absetzung des Bundeskanzlers durch Gerichtsurteil}

Das Grundgesetz setzt offenbar eine gewisse Grundmoral zumindest bei der großen Mehrheit der Parlamentarier voraus.
Sowohl Exekutive als auch Legislative können zwar gegen das Grundgesetz verstoßen, aber sie lehnen die Demokratie und die Gewaltenteilung insgesamt nicht ab und werden sich daher freiwillig vom Bundesverfassungsgericht bei Verstößen korrigieren lassen.

Wenn ein Bundeskanzler Donald Trump spielen will, dann wird ihn das Parlament schon stoppen.
Aber leider muss man davon ausgehen, dass diese Voraussetzung schon bald nicht mehr erfüllt sein könnte.

Sollte tatsächlich erneut ein böser charismatischer Führer wie Hitler auftreten, dem die Massen bereitwillig bis in den Tod folgen würden, könnte ihn kein auch noch so ausgeklügeltes System stoppen.
Er würde jedes System überfluten, weil er überall in jeder Behörde und jeder Gewalt seine Anhänger hätte.

Gegen einen \enquote{törichten Hirten} \footnote{Sacharja 11, 15-17} nach der Art Donald Trumps könnte man aber vielleicht die Demokratie so weit \enquote{härten}, dass sie noch bestehen bliebe, bis er wieder von der Macht verschwände.
Sie würde durch seine Amtszeit zwar Schaden leiden, aber vielleicht ließe sich der unumkehrbare Umbau in eine Autokratie noch verhindern, sodass es noch eine Zeit danach gäbe.

Wenn das Parlament einen außer Kontrolle geratenen Bundeskanzler nicht stoppt, weil zu viele seiner Anhänger darin sitzen, könnte es auch noch ein Gericht tun.
Dazu müsste man aber den Begriff der Immunität von Abgeordneten (und damit des Bundeskanzlers) präziser fassen als bisher.

Man kann Menschen mit allerhand Verdächtigungen und Klagen überhäufen, die sich dann später als grundlos herausstellen.
Allein aber die notwendige Mitwirkung bei den Ermittlungen und die Verteidigung kosten so viel Zeit und Energie, dass sich die Abgeordneten nicht mehr hinreichend um ihre Abgeordnetentätigkeit kümmern können.
Es kann auch sein, dass die Abgeordneten wirklich Vergehen begangen haben, die aber nicht so schwerwiegend waren, dass ihre Sühnung eine Unterbrechung der Abgeordnetentätigkeit rechtfertigen würde.

Um das zu verhindern, ist nach Artikel 46 Absatz 2 und 3 des Grundgesetzes die Genehmigung des Bundestages für Strafverfolgung und jede Einschränkung der persönlichen Freiheit eines Abgeordneten erforderlich.
Diesen Artikel müsste man \enquote{aufbohren} und zusätzlich eine Instanz außerhalb des Bundestages die Abwägungen zwischen Störung der Abgeordneten und notwendiger Strafverfolgung durchführen
lassen, damit eine Strafverfolgung, so sie angemessen und notwendig ist, auch bei Befangenheit des Bundestages durchgeführt werden kann.

\runinhead{Es muss entschieden werden:}

 \begin{enumerate}
     \item Sind einzelne Ermittlungsmaßnahmen, wie z.~B. Hausdurchsuchungen oder Aussagen vor Gericht zulässig?
     (Ermittlungen, die keine störende Wirkung haben und Artikel 47 GG nicht verletzen, sind immer zulässig.)
     Dabei ist zwischen dem öffentlichen Interesse an Aufklärung der Straftat und dem öffentlichen Interesse an einer ungestörten Regierungsarbeit abzuwägen.
     Gleichzeitig muss zwischen der erwarteten Wirksamkeit der Ermittlungsmaßnahme im Hinblick auf die Aufklärung und der durch sie bewirkten Beeinträchtigung des persönlichen Lebens und der Regierungsarbeit abgewogen werden.
     Kann man beispielsweise bei der beantragten Hausdurchsuchung wirklich auf Beweise hoffen, oder ist schon vorher klar, dass sich die Beweise, so es sie gibt, nicht im Haus des Beschuldigten finden lassen werden?
     \item Dürfen Strafen während der Amtszeit oder erst danach vollstreckt werden?
     \item Darf das Amt weiterhin ausgeübt werden?
 \end{enumerate}

Die Abwägungen zu 1 sind klassische Richteraufgaben mit dem Unterschied, dass öffentliches Interesse mit abgewogen werden muss.
Die notwendige Autorität, um eine Ermittlungsmaßnahme gegen den Bundeskanzler abweichend von der Mehrheit des Bundestages zu genehmigen, kann nur dem Bundesverfassungsgericht zukommen, da es die einzige Instanz ist, die auch sonst das Parlament korrigieren kann.
Das Verfahren, um Ermittlungen gegen den Bundeskanzler einzuleiten, wäre also prinzipiell das gleiche wie bei Ermittlungen gegen einfache Parlamentarier: Die zuständige Strafverfolgungsbehörde muss beim Bundestag die Aufhebung der Immunität beantragen.
Wird diese nicht gewährt, können aber einzelne Parlamentarier noch die Aufhebung derselben beim Bundesverfassungsgericht beantragen.
Das Bundesverfassungsgericht muss dann die o.~g. Abwägungen treffen, und die Immunität aufheben oder nicht.
Von sich aus kann das Bundesverfassungsgericht niemals eine Strafverfolgung einleiten oder Immunität aufheben.
Es wird diesbezüglich nur auf Antrag von Parlamentariern tätig.
Man kann noch diskutieren, ob jede einzelne invasive Ermittlungsmaßnahme eine Genehmigung durch das Parlament oder das Bundesverfassungsgericht benötigen sollte, oder ob es nur pauschal eine anlassbezogene Aufhebung der Immunität geben sollte.

2. und 3. hängen zusammen und lassen sich durch einen vom Bundestag frühestens für die nächste Legislaturperiode zu beschließenden Strafenkatalog lösen.
Eine Geldstrafe kann man sicherlich auch während der Amtszeit begleichen, es ist aber klar, dass der Bundeskanzler seine Amtsgeschäfte nicht vom Gefängnis aus führen kann.
Also wenn die Gefängnisstrafe nicht verschoben werden kann, weil das Vergehen so schwer ist, dass seine Sühnung keinen Aufschub duldet, oder weil Vertuschungs- oder Fluchtgefahr besteht, kann der Bundeskanzler nicht im Amt bleiben.
Andersherum, kann es Vergehen geben, die vielleicht nicht mit einer Freiheitsstrafe bewehrt sind, die aber für den Amtsverbleib untragbar sind.
Beide Fälle lassen sich durch einen Katalog von möglichen Vergehen und Strafen abdecken, der Vergehen und/oder Strafen einen Amtsverbleibswert und ggf.~ einen Strafaufschub zuordnet.

In der einfachsten Form könnte der Katalog zum Beispiel vorschreiben, dass die Verhängung von Freiheitsstrafen, die normalerweise nicht zur Bewährung ausgesetzt werden könnten, automatisch eine sofortige Amtsenthebung zur Folge hätte.
Das in Kapitel \ref{sec:abloseregelungen} beschriebene Verfahren würde dann automatisch zur Anwendung kommen.

Der Katalog könnte aber auch explizit Ausnahmen von der Regel nennen: Zum Beispiel, dass Gefängnisstrafen aufgrund von Steuerhinterziehung hinter das Ende der Legislaturperiode verschoben werden können, aber eine erneute Kandidatur nicht erlaubt ist.

Wie sinnvoll eine solche Regelung im Hinblick auf die Vorbildwirkung und die Integrität des Amtes wäre, ist eine andere Frage.
Es wäre jedoch ein Analogon zu der aktuell bestehenden Regelung, dass das Amt vor Strafverfolgung schützt.

In der Verfassung sollte jedenfalls drinstehen, dass Amtsverbleib und Strafaufschub von Regierungsmitgliedern in einem Bundesgesetz geregelt wird, das nicht mit Wirkung auf die aktuelle Legislaturperiode geändert werden darf.

Bei Unstimmigkeiten über die Auslegung dieses Gesetzes entscheidet wie immer das Bundesverfassungsgericht.

\subsection{Parteienverbot und Verwirkung von demokratischen Grundrechten}
Nach Art.~18 GG kann das Bundesverfassungsgericht bei Missbrauch die demokratischen Grundrechte entziehen, darunter auch das passive Wahlrecht.
Nach Art.~21 GG kann das Bundesverfassungsgericht Parteien für verfassungswidrig erklären.
Es wurde vorgeschlagen \footnote{https://www.bundestag.de/resource/blob/658086/50c871fb51049b9edd4153910a0c9398/WD-3-169-19-pdf-data.pdf, Seite 9-10}, für diese Zwecke einen neuen Gerichtshof zu schaffen und das Bundesverfassungsgericht nur die Normenkontrolle durchführen zu lassen.
Diesem Vorschlag schließe ich mich an.
Es sollte möglich sein, bereits gewählten Parlamentsmitgliedern ihr Mandat per Gerichtsbeschluss wieder zu entziehen und ihnen die Wählbarkeit abzuerkennen, wenn sie sich öffentlich
gegen die Demokratie oder die Grundsätze der Verfassung äußern.
Der Wunsch nach Änderung einzelner Vorschriften der Verfassung ist absolut legitim.
Deswegen gibt es ja die Möglichkeit, die Verfassung mit Zweidrittelmehrheit ändern zu können.
Dafür muss man sich auch öffentlich einsetzen dürfen.

Das Problem ist, wenn demokratisch gewählte Menschen die Demokratie praktisch abschaffen oder in eine Ordnung umbauen wollen, die dem freiheitlichen Grundgedanken des Grundgesetzes zuwiderläuft.
Was dieser \enquote{freiheitliche Grundgedanke} des Grundgesetzes ist, muss man genauer präzisieren.
Zurzeit wird er in der Verfassung ja nur durch die sogenannte \enquote{Ewigkeitsklausel} festgehalten.
Man könnte diese ggf. nach Erweiterung und Präzisierung in ein eigenes Dokument, die \enquote{Grundsätze des Grundgesetzes} extrahieren, welches niemals geändert werden darf, solange die Bundesrepublik Deutschland besteht - auch nicht mit Zweidrittelmehrheit.

\subsection[]{Demokratische Notbremse - Kontrolle von Exekutive und Legislative durch den Bundespräsidialrat}\label{subsec:demokratische_notbremse}
\subsubsection{Motivation}\label{subsubsec:demokratische_notbremse_motivation}
Solange das Parlament einwandfrei funktioniert, kann es jede außer Kontrolle geratene Einzelperson inklusive des Bundeskanzlers problemlos stoppen. Leider muss man befürchten, dass dieses vielleicht schon bald nicht mehr der Fall sein könnte. Es reicht, dass eine populistische Partei mit einem Führer nach der Art Donald Trumps die absolute Mehrheit erringt oder im Ergebnis des iterativen Verfahrens zugesprochen bekommt, und schon könnte der Bundeskanzler wie Donald Trump allerhand unvernünftige Dinge tun und der Bundestag unvernünftige Gesetze beschließen, ohne dass irgendjemand sie aufhalten könnte.

Wenn tatsächlich die Mehrheit des Volkes hinter diesen Dingen steht oder zumindest meint, man müsse sie in Kauf nehmen, weil alle anderen Parteien ja versagt haben, dann kann man von innen nicht viel dagegen tun.
Man kann nur hoffen, dass die Menschen aufwachen, bevor es zu spät ist.

Sollte dieses Erwachen bei der Mehrheit des Volkes aber schnell kommen, dann ist es gut, wenn es demokratische Strukturen gibt, die eine vorzeitige Auflösung des Bundestags bei Bedarf erzwingen können.
Es macht keinen Sinn, den Bundestag aufzulösen, wenn wahrscheinlich ist, dass bei Neuwahlen dieselbe populistische Partei erneut gewinnen würde.
Und für den gesellschaftlichen Zusammenhalt und die Demokratie insgesamt wäre es riskant, eine Partei zu verbieten, die von der Mehrheit der Bevölkerung gewünscht wird.
Sollte es aber anders sein, dass die Mehrheit des Volkes seine Meinung schnell ändert, wenn sie am eigenen Leib die Folgen der populistischen Politik zu spüren bekommt, wäre es gut, wenn es ein Instrument gäbe, wodurch eine Regierung auch von außerhalb des Bundestages vorzeitig aus dem Amt gejagt werden könnte.

So ein Instrument zur Einleitung vorgezogener Wahlen von außerhalb des Parlaments wäre auch aus anderen Gründen manchmal wünschenswert.
Auch ohne Populismus kann eine Regierung manchmal so schlecht sein, dass man sie nicht bis zum Ende der Legislaturperiode ertragen möchte.

Die Herausforderung besteht darin, dass man generelle Unzufriedenheit mit der Regierungsarbeit von der Korrektur undemokratischer Entwicklungen unterscheiden muss.
Man kann nicht einfach bei schlechten Umfragewerten automatisch Neuwahlen auslösen.
Sonst kämen wir ja unter Umständen aus dem Wählen nie wieder heraus.

Auch wenn die Unzufriedenheit des Volkes mit der Regierung nur aus der Unfähigkeit der Regierung und nicht aus Populismus herrührt, kann es unter Umständen doch richtig sein,
die Wahlen vorzuziehen, wenn die Unfähigkeit objektiv groß ist und bessere Alternativen in Sicht sind.

Es wird ein nicht dem Populismus unterworfenes, weil nicht vom Bürger gewähltes, Amt gebraucht, was selbst nicht regiert, aber die Regierung in Exekutive und Legislative überwacht.
Es muss die gesamt-gesellschaftliche Lage einschätzen und zum Wohle der Gesamtheit staatspolitische Entscheidungen treffen können.
Dazu gehört unter anderem auch, Abwägungen zu treffen, ob man dem Wunsch des Volkes nach baldigen Neuwahlen stattgeben sollte, oder lieber nicht, weil zu befürchten steht,
dass eine neue Regierung die Sache nicht besser machen würde, oder sowieso wieder dieselben gewählt würden.

Hier kommt der in \ref{subsec:bundespraesidialrat} vorgeschlagene Bundespräsidialrat ins Spiel.
Er soll nicht nur die Demokratie gegen Populismus härten, sondern überhaupt dafür Sorge tragen, dass der Staat als Ganzes funktioniert.
Im Prinzip soll er sich wie ein Bundespräsident bestehend aus sieben Personen verhalten.

\subsubsection[]{Rechte des Bundespräsidialrates} \label{subsubsec:rechte_des_bundespraesidialrats}

\subsubsubsection[]{Vorgezogene Bundestags- oder Kanzlerkandidatennachwahlen auf Volkes Wunsch hin}\label{subsubsubsec:vorgezogene_neuwahlen_auf_volkswunsch_hin}

Der Bundespräsidialrat beobachtet ständig die gesellschaftliche Lage und lässt sich durch das Bundesamt für Meinungsforschung (\ref{subsec:bundesamt_fuer_meinungsforschung}) über die Meinungsverteilung informieren.
Das Amt ermittelt dazu regelmäßig, wie viele Bundesbürger sich zum jetzigen Zeitpunkt vorgezogene Neuwahlen wünschen, und veröffentlicht das Ergebnis.
Wenn eine n-te Mehrheit der Bundesbürger sich vorgezogene Neuwahlen wünscht, kann der Bundespräsidialrat den Bundestag auflösen und Neuwahlen anordnen.
Richtet sich die Unzufriedenheit vor allem gegen den Bundeskanzler, kann er stattdessen auch nur eine Kanzlerkandidatennachwahl anordnen.
Er muss dies aber nicht tun.
Es liegt in seinem Ermessen, dem Volkswillen in dieser Frage nachzugeben oder nicht.

Für n würde ich 3 d.h.~ eine Zweidrittelmehrheit vorschlagen, damit eine Regierung auch noch unpopuläre Entscheidungen durchsetzen kann, was manchmal leider sehr notwendig ist.

\runinhead{Antragsrecht}
Der Bundespräsidialrat kann von sich aus tätig werden.
Außerdem kann jeder einzelne Bundestagsabgeordnete und jede Landesregierung bei ihm beantragen, Neuwahlen anzuordnen, sobald die n-te Mehrheit erreicht ist.
Eine Mindestzahl von Antragstellern sehe ich bewusst nicht vor.
Wenn sich eine populistische Mehrheit gegen das Staatswohl wendet, kommt der Widerspruch aus ihren eigenen Reihen anfangs oft nur von wenigen Abgeordneten.

Lehnt der Bundespräsidialrat einen Antrag ab, muss er dies öffentlich begründen.
Dabei sollte er insbesondere darauf eingehen,
\begin{itemize}
  \item ob eine bessere Alternative in Sicht ist oder bei Neuwahlen voraussichtlich wieder dieselben gewählt würden,
  \item ob die Unzufriedenheit von notwendigen unpopulären Maßnahmen herrührt, deren Wirkung noch aussteht,
  \item wie lange die Legislaturperiode noch dauert und
  \item ob eine akute Krise Neuwahlen zum jetzigen Zeitpunkt unverantwortlich macht.
\end{itemize}
Alle Umfrageergebnisse, auf die er seine Begründung stützt, müssen veröffentlicht werden, damit die Begründung nachprüfbar ist.

\runinhead{Sperrfrist}
Es ist zu erwarten, dass die Opposition jedes Mal einen Antrag stellt, sobald der Neuwahlwunsch über die n-te Mehrheit steigt.
Nach einer Entscheidung des Bundespräsidialrats ist deshalb für X Monate kein neuer Antrag zulässig, gleich von wem.
Ausgenommen sind Anträge, die sich auf eine wesentliche Änderung der Lage stützen, etwa einen deutlichen Anstieg des Neuwahlwunsches oder ein neues Ereignis, das die bisherige Begründung nicht erfasst.
Anträge, die vor seiner Entscheidung eingehen, beantwortet der Bundespräsidialrat gemeinsam.

\runinhead{Keine Wahlsperre}
Anders als bei der Auflösung wegen Staatswohlgefährdung (\ref{absetzung_durch_bundespraesidialrat}) gibt es bei Neuwahlen auf Volkes Wunsch keine Wahlsperre.
Die Abgeordneten haben sich nichts zuschulden kommen lassen, sie haben nur das Vertrauen der Wähler verloren.

\runinhead{Alternative: Pflicht ab einer m-ten Mehrheit}
Man kann diskutieren, ob der Bundespräsidialrat ab einer höheren, m-ten Mehrheit verpflichtet sein sollte, Neuwahlen anzuordnen, wenn diese Mehrheit über Y Monate stabil ist.
Aussetzen dürfte er sie dann nur noch aus eng gefassten Gründen, etwa im Verteidigungsfall oder in einer akuten Krise, und müsste dies öffentlich begründen.
Zusammen bestimmen n und m, wie stark der Volkswille in die Entscheidung für vorgezogene Neuwahlen eingeht: Ab der n-ten Mehrheit darf er berücksichtigt werden, ab der m-ten bindet er.

\runinhead{Gründe dagegen:}
\begin{itemize}
  \item Bisher hat das Volk aus gutem Grund überhaupt keine Möglichkeit, vorgezogene Neuwahlen einzuleiten. Schon eine Möglichkeit, die die Zustimmung des Bundespräsidialrats voraussetzt, ist ein großer Schritt.
  \item Ab der m-ten Mehrheit entscheidet der Bundespräsidialrat nicht mehr, sondern bestätigt nur noch. Populistische Kampagnen hätten ein klares Ziel: den Neuwahlwunsch über die m-te Mehrheit zu treiben.
\end{itemize}

\runinhead{Gründe dafür:}
\begin{itemize}
  \item Die Pflicht schützt davor, dass der Bundespräsidialrat selbst den Volkswillen dauerhaft blockiert, sei es aus Übervorsicht oder weil er unterwandert wurde.
\end{itemize}

Ich würde zunächst auf eine solche Pflicht verzichten.

\subsubsubsection{Ablösung des Bundeskanzlers aufgrund seines Gesundheits- oder Geisteszustandes}\label{subsubsubsec:abloesung_wegen_gesundheit_oder_geisteszustand}
Der Bundespräsidialrat kann eine Untersuchung des Gesundheits- und Geisteszustandes des Bundeskanzlers durch mehrere unabhängige Experten anordnen.
Kommen diese mehrheitlich zu dem Schluss, dass der Bundeskanzler gesundheitlich oder geistig den Anforderungen des Amtes nicht mehr gewachsen ist, kann der Bundespräsidialrat
den Bundeskanzler von seinem Amt entbinden.
Kommen sie einvernehmlich zu diesem Schluss, muss der Bundespräsidialrat den Bundeskanzler von seinem Amt entbinden.

\subsubsubsection{Vetorechte des Bundespräsidialrates}
Der Bundespräsidialrat sollte sein Veto bei Entscheidungen der Bundesregierung oder des Bundestags einlegen dürfen, die das Staatswohl existenziell bedrohen würden.
Das wäre zum Beispiel der Fall,
\begin{enumerate}[label=\alph*)]
    \item wenn der Bundestag zustimmen würde, ein überlegenes Land anzugreifen oder umgekehrt, sich nicht zu verteidigen, wenn man angegriffen wird, weil man aus ideologischen oder sonstigen Gründen, mit dem betreffenden Staat oder seiner Regierung sympathisiert.
     Selbstverständlich hat das Volk das Recht, seinen Staat aufzulösen, wenn es das will und sich in ein anderes Staatengefüge zu integrieren.
     Das bedarf aber eines Volksentscheides und kann nicht per Parlamentsbeschluss geschehen.
    \item Es wäre auch der Fall, wenn eine Entscheidung die Einheit des Staates unmittelbar gefährden würde, indem sie z.B.~ zu Abspaltung von Gebieten oder zu Unruhen oder gar Bürgerkriegen führen würde, oder
    \item wenn eine Entscheidung darauf zielt, die Demokratie zu untergraben, indem z.B.~ die Gewaltenteilung ausgehebelt werden soll.
     Abschaffung der Demokratie zählt als existenziell staatsgefährdend, auch wenn der Fortbestand des deutschen Staates gesichert wäre, da es sich dann im Wesen um einen anderen Staat als den jetzigen handeln würde.
\end{enumerate}

Das Veto des Bundespräsidialrats kann vor dem Bundesverfassungsgericht angefochten werden, indem man z.B. geltend macht, dass die befürchteten existenziellen Folgen gar nicht wirklich eintreten würden.
Der Einspruch kann natürlich keine aufschiebende Wirkung haben, da das Veto ja für existentiell wichtige Entscheidungen gedacht ist, die nicht getroffen und wieder zurückgenommen werden können.

Man kann in Erwägung ziehen, ob der Bundespräsidialrat noch zusätzlich ein weiteres, schwächeres Vetorecht erhalten sollte, welches unter schwächeren Voraussetzungen eingelegt und vom Bundestag mit einer n-ten Mehrheit überstimmt werden kann.
Der Bundespräsidialrat könnte sein schwaches Veto allgemein bei Entscheidungen einlegen, die das Staatswohl oder die Existenz des Staates mittelbar gefährden, wie z.B.~ ohne einen Alternativplan aus der Nato oder der EU auszutreten oder dem Rest der Welt ökonomisch unsinnige Zölle aufzuerlegen.
(Selbstverständlich kann man aus Nato und EU austreten und von jedem Land Zölle erheben, wenn es planvoll im Rahmen einer sinnvollen Risikoabwägung passiert.)

\subsubsubsection{Absetzung des Kanzlers oder Auflösung des Bundestags durch den Bundespräsidialrat}\label{absetzung_durch_bundespraesidialrat}

Der Bundespräsidialrat kann von sich aus sowohl den Bundeskanzler absetzen als auch den Bundestag auflösen, wenn er zu der Einschätzung gelangt, dass diese das Staatswohl gravierend gefährden.
Wenn der Bundespräsidialrat beispielsweise mehrfach sein Veto gegen Entscheidungen einlegen muss, dann kann man davon ausgehen, dass dieser Bundestag bzw. diese Bundesregierung gegen das Staatswohl handeln.
Auch wenn Urteile des Bundesverfassungsgerichts vorsätzlich nicht umgesetzt werden, indem gesetzte Fristen auch nach Ermahnung nicht eingehalten werden, wäre das auch gravierend staatswohlgefährdend und würde auch vorgezogene Neuwahlen rechtfertigen.
Wenn sich eine Pandemie oder andere Seuche ausbreiten würde, und die Bundesregierung würde einfach nichts dagegen unternehmen, wäre das auch gravierend staatswohlgefährdend und Grund, diese zu entlassen.
Vetorechte helfen in diesem Fall nicht, da die Regierung ja eben nichts tut, und man also auch gegen nichts sein Veto einlegen kann.
Der Bundespräsidialrat hat auch kein Recht, konkrete Maßnahmen vorzuschlagen oder die Wirksamkeit der gewählten Maßnahmen zu bewerten.
Das sind Regierungsaufgaben.
Aber einfach nichts tun, wäre ein Grund die Regierung zu feuern.

Ein abgesetzter Bundeskanzler bekommt automatisch eine Sperre und darf innerhalb der nächsten X Jahre nicht mehr kandidieren.
Gleiches gilt für alle Mitglieder der Regierungsfraktion, wenn der Bundestag aufgelöst wird.
Der Bundespräsidialrat kann explizit Ausnahmen von dieser Regel benennen, wenn einzelne Mitglieder durch ihr Stimmverhalten gezeigt haben, dass der Vorwurf der Staatswohlgefährdung auf sie nicht zutrifft.

Die Absetzungs- bzw. Auflösungsentscheidung kann vor dem Bundesverfassungsgericht angefochten werden.
Sie hat u.U. aufschiebende Wirkung.
Das Bundesverfassungsgericht entscheidet über die Dringlichkeit.
In der Hauptsache muss es nicht entscheiden, ob das Staatswohl tatsächlich gravierend gefährdet war, sondern ob man bei gegebener Sachlage davon ausgehen konnte, dass es so war.
Es bleibt ein Ermessensspielraum beim Bundespräsidialrat.
Das Bundesverfassungsgericht stellt aber sicher, dass der Bundespräsidialrat nicht willkürlich Bundestag und Bundeskanzler feuern kann.
Es muss schon eine für alle nachvollziehbare, stichhaltige Begründung vorhanden sein.


\end{document}
